\chapter{Limitations of the Message-Passing Framework}
\label{ch:limitations}
\minitoc

The message-passing framework has been presented in the setting where it
succeeds. This chapter turns to its limits, and two kinds of limit are distinguished. 
The first concerns propagation --- what happens to representations as messages are 
passed and when the homophily assumption behind aggregation no longer holds. 
The second concerns expressive power --- what the 
framework can distinguish at the graph level, independent of training.

\section{Pathologies of Neighborhood Propagation}
\label{sec:propagation-pathologies}

This section examines three ways in which neighborhood propagation can degrade the
representations it produces: \textit{oversmoothing}, the collapse of
representations under depth; \textit{oversquashing}, the distortion of information
drawn from far away; and \textit{heterophily}, where neighboring nodes tend to
be dissimilar rather than similar.

\subsection{Oversmoothing}
\label{sec:oversmoothing}

A GCN layer smooths, reducing variation across edges and making representations 
of neighboring nodes more similar (Section~\ref{sec:gcn}). A few layers may be sufficient 
to align the representations with the labels on a homophilous graph; the difficulty is that 
smoothing does not stop at a useful point. Applied repeatedly, the operation that makes 
adjacent nodes more similar makes all nodes similar: as depth grows, the representations 
within each connected component converge toward a common vector. 
This is \textit{oversmoothing}, first identified for the GCN 
as a consequence of its Laplacian smoothing~\cite[11.5pc]{li2018}.

\begin{figure}
  \centering
  \includegraphics{oversmoothing}
  \caption[Node representations collapse as depth increases.]{%
    Node representations at increasing depth. With few layers the representations
    of different nodes remain distinct; as layers are stacked, repeated smoothing
    drives them together, until at large depth the nodes of each connected
    component converge toward a common representation and can no longer be told apart.
  }
  \label{fig:oversmoothing}
\end{figure}

This places a ceiling on useful depth that is in direct tension with the receptive
field of Section~\ref{sec:receptive-field}. The useful depth of a network is therefore 
a compromise between the receptive field and oversmoothing. 
Architectural designs that relieve this tension are revisited in Chapter~\ref{ch:conclusions}.

\subsection{Oversquashing}
\label{sec:oversquashing}

The receptive field grows with depth, suggesting that a deep enough
network lets any node draw on any other. However, the number of
nodes within $k$ hops typically grows exponentially with $k$, yet the
representation summarizing this region is a vector of fixed dimension, the same
size whether it gathers a handful of nodes or thousands. Information from an
exponentially growing region is compressed into a fixed-size representation,
making information loss unavoidable.

\begin{figure}[h]
  \centering
  \includegraphics{oversquashing}
  \caption[A growing receptive field squashed into a fixed-size vector.]{%
    The receptive field of a node, unrolled over layers, typically grows exponentially: each
    additional hop multiplies the number of nodes that must be summarized. The
    representation gathering this region is a vector of fixed size, independent of
    how many nodes feed into it, so information from the wider neighborhood is
    compressed, reducing the contribution of any single distant node.
  }
  \label{fig:oversquashing}
\end{figure}

This is \textit{oversquashing} (Figure~\ref{fig:oversquashing}): the distortion that results when information from a
node's wider neighborhood is squeezed through the fixed-size representations of the
intermediate nodes~\cite[-4.25pc]{alon2021}. Its consequence is that message-passing
networks struggle on tasks depending on long-range interaction. To reach a distant
target, a signal must pass through every intermediate representation, and at each
step it is aggregated with everything else in that node's neighborhood.
The framework can reach distant nodes by adding depth, but the 
information it gathers from them can arrive too degraded to be useful.

Oversquashing and oversmoothing both worsen with depth but are distinct failures.
Oversmoothing is a loss of distinction \emph{between nodes}: their representations
converge. Oversquashing is a loss of information \emph{from distance}: one node's
representation cannot faithfully hold its large receptive field. Together they
bound how deep a network can usefully be made.

\subsection{Heterophily}
\label{sec:heterophily}

Oversmoothing and oversquashing are costs of the mechanism that arise whatever the
data. Heterophily is instead a failure of fit between mechanism and data.
A GCN layer rests on the assumption that
adjacent nodes are similar --- the homophily of
Section~\ref{sec:structural-properties}.
When it holds, aggregating a neighborhood
reinforces what a node indicates; when it fails, the same aggregation works against
the task.

The degree to which a graph is homophilous can be measured, for a
node-classification task, by the fraction of edges joining nodes of the same class.
Writing $y_i$ for the class label of node $i$, the \textit{edge homophily} of a
graph is
\begin{equation}
  h(\mathcal{G})
    = \frac{\bigl\lvert
        \{\, \{i, j\} \in E : y_i = y_j \,\}
      \bigr\rvert}{m},
  \label{eq:edge-homophily}
\end{equation}
the number of edges whose endpoints share a label divided by the total number of
edges $m$. A value near one describes a strongly homophilous graph; a value near
zero, a strongly heterophilous one, in which edges predominantly join nodes of
different classes.

On a heterophilous graph, the aggregation misleads: a node's neighbors mostly
belong to other classes, so smoothing their representations pulls the node toward
classes it does not belong to, blurring the distinction the task requires. 
The framework does not fail to use the graph;
it uses the graph through an assumption the graph does not satisfy. Aggregation
schemes that separate a node's own contribution from its neighbors', or learn to
weight neighbors unequally, relax this assumption and are noted in
Chapter~\ref{ch:conclusions}.

\section{The Expressivity Ceiling under 1-WL}
\label{sec:expressivity-ceiling}

Section~\ref{sec:expressive-power} established that the expressive power of message
passing is bounded above by the \mbox{1-WL} test, and that the architecture of
Section~\ref{sec:gin} was designed to reach this ceiling. 
This section makes the limit concrete by exhibiting what it excludes.

What the ceiling excludes is a pair of non-isomorphic graphs that \mbox{1-WL} ---
and therefore every message-passing network --- colors identically. The simplest
examples arise among pairs of \textit{regular} graphs, in which every node has the
same degree. All nodes begin with the same
color; and whenever every node shares a color, each one has the same number of
identically colored neighbors, so the refinement assigns them all the same new color.
The condition therefore reproduces itself at every round, and the test never
separates the nodes, so two such graphs receive the same final coloring. Yet regular
graphs of the same degree need not be isomorphic, differing in structural properties
the count of same-colored neighbors never exposes, such as whether they form one
connected cycle or two disjoint shorter ones (Figure~\ref{fig:wl-indistinguishable}).

\begin{figure}
  \centering
  \includegraphics{wl-indistinguishable}
  \caption[Two graphs that 1-WL cannot distinguish.]{%
    Two non-isomorphic regular graphs in which every node has the same degree. The
    \mbox{1-WL} test colors all nodes of both graphs identically and never separates
    them, so no message-passing network can tell the two graphs apart, although they
    are not isomorphic.
  }
  \label{fig:wl-indistinguishable}
\end{figure}

This is a genuine ceiling, not a shortcoming of present architectures, because it
follows from the locality of the computation: a network that only ever sees
aggregated neighborhoods cannot distinguish nodes whose neighborhoods are
themselves indistinguishable. Surpassing it means
changing what the computation may see. Such methods
fall outside the present framework; they are taken up as future work
in Chapter~\ref{ch:conclusions}.